% Einheit 2 von 4 – Rückwärts (bank/kenngroessen-von-verteilungen/e2.jsonl, zusammenbau v0.1)
%% TODO Zeitmarke Einheit 2: Marken-Zeile nicht lesbar
%% TODO Zweigzeile Teil 1: was hier gelernt wird (Satz in Schülersprache aus dem Einheitstitel) – Einheit 2
\einheitenkopf[e2]{Einheit 2 von 4 · Rückwärts}
\zweigzeile{Abitur GK}
\setcounter{aufgabe}{13}

%% TODO Titel: Ich-kann-Satz fehlt in der Bank (Platzhalter: Kettenname) – Nr. 14
%% TODO Anweisung: ein Satz über der ersten Teilaufgabe fehlt in der Bank – Nr. 14
\begin{aufgabe}{Rückwärts}
%% TODO \rechenplatz{n}: Schrittzahl des Grundfalls fehlt in der Bank (nur bei fester Schreibform, 2.3 b)
\begin{teile}
\teil Ein Würfelspiel zahlt bei einer Sechs 4\,€ aus, sonst nichts. Gefragt ist, wie viel im Mittel je Spiel ausgezahlt wird. Vorwärts (Erwartungswert berechnen) oder rückwärts (Unbekannte aus einer Erwartungswertbedingung)? \\ \kreuz{vorwärts} \\ \kreuz{rückwärts}
\end{teile}
\begin{gleichungsraster}[2]
\gl{\text{Ein Spiel kostet 2\,€ Einsatz. Mit der Wahrscheinlichkeit 0{,}25 wird der Betrag $a$ (in Euro) ausgezahlt, sonst nichts. Für welches $a$ ist das Spiel fair?}} &
\gl{\text{Ein Würfel trägt auf vier Seiten die Zahl 1 und auf zwei Seiten die Zahl $b$. $X$ ist die geworfene Zahl, und es gilt $E(X) = 3$. Wie groß ist $b$?}} \\
\gl{\text{In einer Urne liegen sechs weiße und $x$ schwarze Kugeln. Zwei Kugeln werden mit einem Griff gezogen; gewonnen wird, wenn beide weiß sind. Wie viele schwarze Kugeln müssen in der Urne liegen, damit die Gewinnwahrscheinlichkeit $\frac{1}{3}$ beträgt (Abitur 2018 GK)?}} &
\gl{\text{Eine Basketballspielerin trifft bei jedem Freiwurf mit der Wahrscheinlichkeit 0{,}7. Wie viele Freiwürfe muss sie mindestens werfen, damit sie im Mittel mehr als 30 Treffer erzielt?}} \\
\end{gleichungsraster}
\begin{teile}
\teil Ein Würfel hat die sichtbaren Seiten 2, 2 und 6. Die drei verdeckten Seiten sollen mit Zahlen aus 1, 2, 3, 4 beschriftet werden (Wiederholung erlaubt). Bedingungen: der Erwartungswert beim einmaligen Werfen ist 3; auf dem Würfel stehen genau drei verschiedene Zahlen; die Wahrscheinlichkeit, bei zwei Würfen zweimal dieselbe Zahl zu werfen, ist $\frac{7}{18}$. Untersuche, ob es eine Beschriftung mit allen drei Eigenschaften gibt \hfill (Abitur 2024 LK).
\end{teile}
\end{aufgabe}

%% TODO Titel: Ich-kann-Satz fehlt in der Bank (Platzhalter: Kettenname) – Nr. 15
%% TODO Anweisung: ein Satz über der ersten Teilaufgabe fehlt in der Bank – Nr. 15
\begin{aufgabe}{Erwartungswertgleichung für einen Glücksradparameter aus den Spielregeln herleiten}
\begin{gleichungsraster}[2]
\gl{\text{Ein Glücksrad hat Felder mit den Zahlen 4 und 1; die 4 erscheint mit der Wahrscheinlichkeit $q$. Es wird zweimal gedreht; der Rabatt in Prozent ist das Produkt der beiden Zahlen. Auf lange Sicht soll der Rabatt im Mittel 4\,\% betragen. Zeige, dass $q$ die Gleichung $3q^2 + 2q - 1 = 0$ erfüllt, und bestimme $q$ (Abitur 2024 GK).}} & \\
\end{gleichungsraster}
\end{aufgabe}

%% TODO Titel: Ich-kann-Satz fehlt in der Bank (Platzhalter: Kettenname) – Nr. 16
%% TODO Anweisung: ein Satz über der ersten Teilaufgabe fehlt in der Bank – Nr. 16
\begin{aufgabe}{Zwei Wahrscheinlichkeiten einer Verteilung aus dem Erwartungswert und der Summe 1 bestimmen}
\begin{gleichungsraster}[2]
\gl{\text{Eine App vergibt beim täglichen Öffnen 5, 10 oder 40 Münzen; 40 Münzen gibt es weiter mit der Wahrscheinlichkeit 0{,}05. Die Wahrscheinlichkeiten für 5 und 10 Münzen werden so geändert, dass in 30 Tagen im Mittel 270 Münzen anfallen. Bestimme die beiden neuen Wahrscheinlichkeiten (Abitur 2021 GK).}} & \\
\end{gleichungsraster}
\end{aufgabe}

%% TODO Titel: Ich-kann-Satz fehlt in der Bank (Platzhalter: Kettenname) – Nr. 17
%% TODO Anweisung: ein Satz über der ersten Teilaufgabe fehlt in der Bank – Nr. 17
\begin{aufgabe}{Restwahrscheinlichkeit und Erwartungswert eines Teilgewinns aus dem Gesamterwartungswert bestimmen}
\begin{teile}
\teil Bei einem Gewinnspiel gewinnt jeder Teilnehmer: einen Gutschein über 50\,€ mit der Wahrscheinlichkeit 0{,}002, einen über 10\,€ mit 0{,}01, sonst einen Sachpreis. Im Mittel gewinnt ein Teilnehmer Preise im Wert von 95 Cent. Bestimme die Wahrscheinlichkeit für einen Sachpreis und den Erwartungswert des Gewinns, der auf die Sachpreise entfällt \hfill (Abitur 2023 LK). $P = $ \leerfeld , \leerfeld[€]
\end{teile}
\end{aufgabe}

%% TODO Titel: Ich-kann-Satz fehlt in der Bank (Platzhalter: Kettenname) – Nr. 18
%% TODO Anweisung: ein Satz über der ersten Teilaufgabe fehlt in der Bank – Nr. 18
\begin{aufgabe}{Verhältnis zweier Auszahlungen aus dem Ausgleich der Erwartungswerte berechnen}
\begin{teile}
\teil Bei Spiel A wird mit der Wahrscheinlichkeit 0{,}3 der Betrag $a$ ausgezahlt, bei Spiel B mit der Wahrscheinlichkeit 0{,}2 der Betrag $b$, sonst jeweils nichts. Beide Spiele sollen den gleichen Erwartungswert der Auszahlung haben. In welchem Verhältnis stehen $a$ und $b$? $a : b = $ \leerfeld
\end{teile}
\end{aufgabe}

%% TODO Titel: Ich-kann-Satz fehlt in der Bank (Platzhalter: Kettenname) – Nr. 19
%% TODO Anweisung: ein Satz über der ersten Teilaufgabe fehlt in der Bank – Nr. 19
\begin{aufgabe}{Kugelzahl für den größten Erwartungswert der Auszahlung ermitteln}
\begin{teile}
\teil In einer Urne liegen zehn Kugeln, davon $k$ rote, die übrigen weiß. Es wird zweimal mit Zurücklegen gezogen; bei zwei verschiedenen Farben gibt es 1\,€, sonst nichts. Für welches $k$ ist der Erwartungswert der Auszahlung am größten, und wie groß ist er dann? $k = $ \leerfeld , $E = $ \leerfeld[€]
\end{teile}
\end{aufgabe}

%% TODO Titel: Ich-kann-Satz fehlt in der Bank (Platzhalter: Kettenname) – Nr. 20
%% TODO Anweisung: ein Satz über der ersten Teilaufgabe fehlt in der Bank – Nr. 20
\begin{aufgabe}{Untere Schranke für eine Kugelbeschriftung aus dem Erwartungswert der Summe ohne Rechnung begründen}
\begin{teile}
\teil In einer Urne liegen drei Kugeln mit den Zahlen 1, 2 und $c$. Zwei Kugeln werden ohne Zurücklegen gezogen; der Erwartungswert ihrer Summe ist 6. Begründe ohne Rechnung, dass $c$ größer als 3 sein muss.
\end{teile}
\end{aufgabe}

%% TODO Titel: Ich-kann-Satz fehlt in der Bank (Platzhalter: Kettenname) – Nr. 21
%% TODO Anweisung: ein Satz über der ersten Teilaufgabe fehlt in der Bank – Nr. 21
\begin{aufgabe}{Aussage über gleiche Erwartungswerte aufeinanderfolgender Parameterwerte über eine Gleichung beurteilen}
\begin{teile}
\teil In einer Urne liegen neun Kugeln, $n$ rote und $9 - n$ weiße. Es wird zweimal mit Zurücklegen gezogen; bei zwei verschiedenen Farben gibt es 27\,€. Der Erwartungswert der Auszahlung ist $\frac{2n(9 - n)}{3}$ Euro. Beurteile mit einer Gleichung die Aussage: „Es gibt zwei aufeinanderfolgende Werte von $n$ mit gleichem Erwartungswert.“
\end{teile}
\end{aufgabe}

%% TODO Titel: Ich-kann-Satz fehlt in der Bank (Platzhalter: Kettenname) – Nr. 22
%% TODO Anweisung: ein Satz über der ersten Teilaufgabe fehlt in der Bank – Nr. 22
\begin{aufgabe}{Rückwärts – Fehler finden, Begründen, Anwendung}
\begin{teile}
\teil Ein Spiel kostet 2\,€ Einsatz; mit der Wahrscheinlichkeit 0{,}25 wird der Betrag $a$ ausgezahlt, sonst nichts. Tim sucht $a$ für ein faires Spiel. Finde den Fehler und rechne richtig. \rechnung{0{,}25 \cdot (a - 2) + 0{,}75 \cdot (-2) &= 2 \\ 0{,}25a - 2 &= 2 \\ a &= 16}
\teil Begründe, warum die Angabe „auf lange Sicht gleichen sich Einsätze und Auszahlungen aus“ eine Gleichung für eine unbekannte Auszahlung liefert.
\teil Eine Handyversicherung zahlt bei einem Schaden im Mittel 250\,€; 8\,\% der versicherten Geräte haben im Jahr einen Schaden. Je Vertrag entstehen zusätzlich 12\,€ Verwaltungskosten. Wie hoch muss die Jahresprämie mindestens sein, damit die Versicherung im Mittel keinen Verlust macht? \leerfeld[€]
\end{teile}
\end{aufgabe}
